
et on
définit un homomorphisme :
@@M0@@
en posant :
@@M1@@
avec @@M2@@.

\textbf{Proposition 4.7.5.} \uline{Les opérations définies par} (4.7.1) \uline{et} (4.7.2) \uline{munissent} @@M3@@ \uline{d'une structure de} @@M4@@-\uline{anneau}, \uline{et} @@M5@@ \uline{est un homomorphisme injectif qui identifie} @@M6@@, \uline{muni des opérations définies par} (4.7.3) \uline{et} (4.7.4) \uline{à un} @@M7@@-\uline{idéal de} @@M8@@.

Il est immédiat de vérifier que @@M9@@ est un @@M10@@-anneau : on se réfèrera par exemple à V 5.5, où une construction tout à fait analogue est faite. La démonstration de la seconde partie de la proposition se fait également sans difficultés en utilisant les relations précédentes, la relation @@M11@@ qui résulte immédiatement des hypothèses, la relation @@M12@@ (voir VI 5.9) et les relations de V 5.

On pose alors :
@@M13@@
et on note @@M14@@ l'homomorphisme canonique @@M15@@. On considère @@M16@@ comme @@M17@@-anneau grâce à la structure quotient.

On définit ensuite des homomorphismes additifs :
@@M18@@

L'homomorphisme @@M19@@ est le composé @@M20@@. De même, l'homomorphisme @@M21@@ est le composé @@M22@@. On a alors :
@@M23@@
